\section{Purpose: reconstruct admissible magnetics from electrical samples}
An accurate interpolator can match two measured flux surfaces without making
them the gradient of any common magnetic potential. If those surfaces are
later differentiated, visually small inconsistencies may become large errors
in differential inductance. The research question here is therefore not only
how to minimize pointwise flux error, but how to reconstruct smooth saturated
$dq$ maps while retaining known magnetic structure.
An additional question is which components of such a map are independently
observable through electromagnetic torque, and how this information
complements electrical observations and differential consistency.

We begin with terminal voltage, current, and electrical speed. The stationary
voltage balance yields a reconstructed terminal-current map. To interpret
its vectors as slopes of a latent scalar co-energy, one must additionally
establish a conservative storage description in conjugate magnetic current
coordinates. Loss currents, measurement bias, and state reduction can prevent
that identification. The proposed potential models the storage subsystem;
it is not an unrestricted model of the entire measured terminal behavior.
Reflection symmetry constrains that surface, its gradient supplies
both flux components, and its Hessian supplies their differential inductances.
Energy-based machine descriptions and saturation models support this approach
\cite{jebai2014,jebai2011}; standard voltage conventions follow
\cite{binder2017,schroeder2021}.

This paper owns the full co-energy, symmetry, integrability, and fitting
development. Its companion studies the complementary question: how do
departures from this admissible structure constrain interpretations of
resistance, rotor-angle, and related reconstruction errors
\cite{companion-diagnostics}? Here those
errors motivate preserving fit residuals; they are not estimated. Both papers
remain independently readable and carry their own framework and numerical
data. The gradient-fitting study concerns noise, sampling, and derivative
quality. A separate torque study resolves observable directions and
tests remaining ambiguities with known synthetic fields.
An additional real-data admissibility study confronts the magnetic class with
terminal-current maps. Exact integration reproduces a disagreement between
symmetry and circulation equivalents; reducing high-current parity residuals
does not improve global consistency. This study tests assumptions before
fitting and does not claim successful real-data potential recovery.
The torque development is placed between differential inductance and fitting:
it shows why even a reciprocal map has radial components invisible to torque,
then adds the torque row to a common observation formulation.
A multi-speed real comparison follows a signal-provenance gate. Its CAN
channel is not qualified as electromagnetic torque, so the numerical
disagreement is reported conditionally, without a real correction claim.

Unadorned quantities denote a true constitutive field, measured electrical
inputs are explicitly marked, $\flux^{\rm rec}$ denotes reconstructed samples,
and $\hat\flux_{\rm fit}=\nabla\hat\coenergy$ denotes a fitted field. All magnetic
statements refer to fixed temperature and magnet state, with an effective
fundamental or consistently averaged quasi-static description. No additional
energy measurement is assumed.
\begin{figure}[tb]
\centering
\begin{tikzpicture}[node distance=8mm,]
 \node[pipeline block] (m) {Measured $u_d,u_q,i_d,i_q,\omega_e$\\and used $R_s$};
 \node[pipeline block,below=of m] (f) {Stationary reconstruction\\$\vect\psi^{\rm rec}(i_d,i_q)$ samples};
 \node[pipeline block,below=of f] (energyfit) {Declared storage coordinates and assumptions\\Symmetric potential gradient fit};
 \node[pipeline block,below=of energyfit] (d) {Analytic derivatives\\$\hat\flux_{\rm fit}=\nabla\hat W'$, $\hat\Ldiff=\nabla^2\hat W'$};
 \node[diagnostic block,text width=39mm,right=12mm of energyfit] (r)
  {Consistency channels\\Electrical observations\\Reciprocal derivatives\\Qualified torque projection\\[2pt]
   Preserve structured residuals};
 \draw[direction arrow] (m)--(f);\draw[direction arrow] (f)--(energyfit);
 \draw[direction arrow] (energyfit)--(d);\draw[direction arrow] (energyfit)--(r);
\end{tikzpicture}
\caption{Measurement-to-model pipeline. No energy sensor is required:
the gradient fit defines a storage model under declared coordinate and
admissibility assumptions. Inconsistent terminal samples yield a model
projection. Electrical, differential, and mechanical consistency have
different information and nullspaces; their unexplained residuals remain
available for the companion's error hypotheses.}
\label{fig:pipeline}
\end{figure}

